% Header: Here are all packages used and some additional definitions
%%%%%%%%%%%%%%%%%%%%%%%%%%%%%%%%%%%%%%%%%%%%%%%%%%%%%%%%%%%%%%%%%%%

\documentclass[11pt,a4paper]{article}
\usepackage[margin=2.5cm]{geometry}
\usepackage[onehalfspacing]{setspace}
\usepackage{graphicx} % zum Einbinden von Graphiken
\usepackage[breaklinks=true,colorlinks=true,linkcolor=blue,urlcolor=blue,citecolor=blue]{hyperref} % f. Referenzen
\usepackage{amsmath,amsthm,amssymb} % Mathematik Umgebung 
\usepackage{icomma} % Intelligentes Komma, das den richtigen Abstand zwischen Dezimalzahlen als auch in Formeln wählt.
%\usepackage[ngerman]{babel} % Deutsche Bezeichnungen bei Inhaltsangabe etc
\usepackage[T1]{fontenc}    % andere Schriftsatzkodierung für richtige Silbentrennung bei Umlauten
\usepackage[locale = DE,space-before-unit=true,per-mode = symbol]{siunitx} % Bessere Einheiten
\usepackage{placeins} % Definiert den Befehl “\FloatBarrier”, der die Ausgabe der davor eingebundenen Bilder erzwingt, befor der Text weiter geht. (Mit vorsicht zu verwenden)


\graphicspath{{Bilder/}}

\DeclareSIUnit{\dBm}{dBm}
\DeclareSIUnit[per-mode=reciprocal]\WN{\per\centi\meter}

\usepackage{xcolor}
\usepackage{tabularray}
\newcommand{\up}{\ensuremath{\uparrow}}
\newcommand{\down}{\ensuremath{\downarrow}}

% --- Scaffolding helpers (to be filled in paper-by-paper) ---
% Placeholder for content still to be completed from the source paper.
\newcommand{\tbd}{\textcolor{gray}{\emph{[TBD]}}}
% Stacked reference cell: cohort name + linked citation, BioProject ID underneath.
% Usage: \refcell{CohortName}{bibkey}{PRJ...}
\newcommand{\refcell}[3]{\shortstack[l]{#1~\cite{#2}\\[2pt]{\scriptsize #3}}}


\usepackage[backend=biber, style=numeric, sorting=none, autocite = superscript, maxnames = 10]{biblatex}
\addbibresource{references.bib}

\AtBeginBibliography{\renewbibmacro*{begentry}{\printnames{author}}}
\AtEveryBibitem{\clearfield{note}}

\setlength{\bibhang}{1em}
\defbibenvironment{bibliography}
  {\list
     {\printtext[labelnumberwidth]{%
        \printfield{labelprefix}%
        \printfield{labelnumber}}}
     {\setlength{\labelwidth}{\labelnumberwidth}%
     \setlength{\leftmargin}{\bibhang}
     % \setlength{\leftmargin}{\labelwidth}%
      \setlength{\labelsep}{\biblabelsep}%
      \setlength{\itemsep}{\bibitemsep}%
      \setlength{\parsep}{\bibparsep}%
      \addtolength{\leftmargin}{\labelsep}%
      \setlength{\itemindent}{-\bibhang}
      %\setlength{\itemindent}{0pt}%
      \renewcommand*{\makelabel}[1]{\hss##1}%
      \raggedright}}
  {\endlist}
  {\item}


%%%%%%%%%%%%%%%%%%%%%%%%%%%%%%%%%%%%%%%%%%%%%%%%%%%%%%%%%%%%%%%%%%%
\begin{document}
%
\title{Gut Microbiome in Adult Allo-SCT Patients: Analysis Sub Cohort}
\author{Sophie Huebler \thanks{\href{mailto:sophie.huebler@utah.edu}{sophie.huebler@utah.edu}}}
\date{\today}
\maketitle
%\vfill
\renewcommand\abstractname{Objective}
\section*{\abstractname}
\textit{Make a list of adult allo-HSCT recipient cohorts with publicly available microbiome data (and associated clinical data).}
\thispagestyle{empty}
%
%
\tableofcontents
\thispagestyle{empty}
\cleardoublepage
\pagenumbering{arabic} 
\newpage
%
%
\section{Studies}

% =====================================================================
% TABLE 1 — Study-level summary
% Columns: Year | Patients | Study design | Main results | Reference
% Reference cell (rightmost) stacks: cohort name + \cite, BioProject underneath.
% Data cells left as \tbd — fill in paper-by-paper.
% =====================================================================
\begin{table}[h]
\centering
\caption{Human studies of the association between the gut microbiota and GVHD with
publicly available 16S sequencing and clinical metadata included in the analytic sub-cohort.}
\begin{tblr}{
  width = \linewidth,
  colspec = {Q[l,m,0.07\linewidth] Q[l,m,0.12\linewidth] Q[l,m,0.16\linewidth] Q[l,m,0.42\linewidth] Q[l,m,0.17\linewidth]},
  row{1} = {font=\bfseries}, % Bold header row
  hline{1,2,Z} = {0.8pt},    % Rules at top, below header, and bottom
  cells = {font=\small},
  column{5} = {font=\footnotesize}, % Reference column smaller
}

Year & Patients* & Study design & Main results & Reference \\
\tbd & \tbd & \tbd & \tbd & \refcell{Liu2017}{Liu2017}{PRJEB16057} \\
\tbd & \tbd & \tbd & \tbd & \refcell{Fujimoto2024}{Fujimoto2024}{PRJNA1109929} \\
\tbd & \tbd & \tbd & \tbd & \refcell{Artacho2024}{artacho2024multimodal-9e5}{PRJNA1088414} \\
\tbd & \tbd & \tbd & \tbd & \refcell{Ingham2019}{ingham2019specific-2d2}{PRJEB25221} \\
\tbd & \tbd & \tbd & \tbd & \refcell{Vallet2023}{Vallet2023}{PRJNA902819} \\ % TODO: add Vallet2023 to references.bib
\tbd & \tbd & \tbd & \tbd & \refcell{DAmico2019}{DAmico2019}{PRJNA592853} \\ % TODO: add DAmico2019 to references.bib
\tbd & \tbd & \tbd & \tbd & \refcell{Jarosch2023}{jarosch2023multimodal-0f7}{PRJEB60178} \\
\end{tblr}
\vspace{2pt}
\begin{flushleft}
\footnotesize \textit{*Patients refers to the number of patients with 16S microbiome
measurements included here. Main results give a brief description of the primary
findings of the original study.}
\end{flushleft}
\end{table}

% =====================================================================
% TABLE 2 — Sample generation / sequencing characteristics
% Columns: Reference | Platform | Region | Primers | Reads | Avg. depth
% (Platform = Illumina/454 pyro/etc.; Region = V4, V3-V4, etc.;
%  Reads = paired / single-end; Avg. depth = mean reads per sample.)
% Data cells left as \tbd — fill in paper-by-paper.
% =====================================================================
\begin{table}[h]
\centering
\caption{Sample generation and 16S sequencing characteristics for each cohort.}
\begin{tblr}{
  width = \linewidth,
  colspec = {Q[l,m,0.16\linewidth] Q[l,m,0.14\linewidth] Q[l,m,0.10\linewidth] Q[l,m,0.22\linewidth] Q[l,m,0.12\linewidth] Q[l,m,0.12\linewidth]},
  row{1} = {font=\bfseries},
  hline{1,2,Z} = {0.8pt},
  cells = {font=\small},
  column{1} = {font=\footnotesize},
}

Reference & Platform & Region & Primers & Reads & Avg. depth \\
Liu2017 & \tbd & \tbd & \tbd & \tbd & \tbd \\
Fujimoto2024 & \tbd & \tbd & \tbd & \tbd & \tbd \\
Artacho2024 & \tbd & \tbd & \tbd & \tbd & \tbd \\
Ingham2019 & \tbd & \tbd & \tbd & \tbd & \tbd \\
Vallet2023 & \tbd & \tbd & \tbd & \tbd & \tbd \\
DAmico2019 & \tbd & \tbd & \tbd & \tbd & \tbd \\
Jarosch2023 & \tbd & \tbd & \tbd & \tbd & \tbd \\
\end{tblr}
\vspace{2pt}
\begin{flushleft}
\footnotesize \textit{Platform: sequencing platform (e.g.\ Illumina MiSeq, 454
pyrosequencing). Region: 16S hypervariable region targeted (e.g.\ V4, V3--V4).
Reads: paired-end or single-end. Avg.\ depth: mean reads per sample after
demultiplexing.}
\end{flushleft}
\end{table}

\FloatBarrier
\newpage

% =====================================================================
% Per-study sections. Each has three subsections:
%   1. Summary            — short summary of the paper
%   2. Identified species — microbial species discussed in the original paper
%   3. Processing          — pipeline steps applied, primer sequences, layout, etc.
% Fill in paper-by-paper.
% =====================================================================

\section{Liu2017}
\subsection{Summary}
\tbd

\subsection{Identified microbial species}
\tbd

\subsection{Processing}
% Pipeline steps hit (single/paired-end, forward-read-only rescue, etc.),
% primer sequences, trimming/DADA2 parameters, reference DB.
\tbd

\section{Fujimoto2024}
\subsection{Summary}
\tbd

\subsection{Identified microbial species}
\tbd

\subsection{Processing}
\tbd

\section{Artacho2024}
\subsection{Summary}
\tbd

\subsection{Identified microbial species}
\tbd

\subsection{Processing}
\tbd

\section{Ingham2019}
\subsection{Summary}
\tbd

\subsection{Identified microbial species}
\tbd

\subsection{Processing}
\tbd

\section{Vallet2023}
\subsection{Summary}
\tbd

\subsection{Identified microbial species}
\tbd

\subsection{Processing}
\tbd

\section{DAmico2019}
\subsection{Summary}
\tbd

\subsection{Identified microbial species}
\tbd

\subsection{Processing}
\tbd

\section{Jarosch2023}
\subsection{Summary}
\tbd

\subsection{Identified microbial species}
\tbd

\subsection{Processing}
\tbd


\newpage
\section{Bibliography}
\printbibliography[heading=bibintoc]


\end{document}